% Required preamble commands/packages:
% \usepackage{amsmath,amssymb,amsthm}
% \newtheorem{theorem}{Theorem}
% \newcommand{\E}{\mathbb{E}}
% \newcommand{\R}{\mathbb{R}}
% \newcommand{\diag}{\operatorname{diag}}
% \newcommand{\argmin}{\operatorname*{arg\,min}}

\section{Full Proof of Theorem 1}
\label{app:theorem1_full_proof}

\begin{theorem}[Slow subspace identifiability]
Consider the stable linear-Gaussian dynamics
\begin{equation}
    x_{t+1}=Ax_t+w_t,
    \qquad
    w_t \sim \mathcal{N}(0,Q),
\end{equation}
with stationary covariance $\Sigma_x=\E[x_t x_t^\top]\succ 0$. Assume that
$A$ is diagonalizable and admits a slow-fast spectral gap such that the $r$
slow modes have strictly smaller derivative variance than the remaining modes.
Then any population minimizer of
\begin{equation}
    \min_{U \in \R^{d \times r}}
    \E\left[\|U^\top(x_{t+1}-x_t)\|_2^2\right]
    \quad
    \text{subject to}
    \quad
    U^\top \Sigma_x U = I_r
\end{equation}
spans the true slow subspace up to rotation.
\end{theorem}

\begin{proof}
Let
\begin{equation}
    \dot{x}_t = x_{t+1}-x_t.
\end{equation}
In discrete time, this is the finite-difference version of the temporal
derivative used by the slow-feature objective. Substituting the dynamics gives
\begin{equation}
    \dot{x}_t = (A-I)x_t+w_t.
\end{equation}
The derivative covariance is therefore
\begin{align}
    \Sigma_{\dot{x}}
    &=
    \E[\dot{x}_t\dot{x}_t^\top] \\
    &=
    (A-I)\Sigma_x(A-I)^\top
    + Q
    + \E[(A-I)x_t w_t^\top]
    + \E[w_t x_t^\top(A-I)^\top].
\end{align}
Under the usual state-space assumption that $w_t$ is independent of $x_t$, the
cross terms vanish, so
\begin{equation}
    \Sigma_{\dot{x}} = (A-I)\Sigma_x(A-I)^\top + Q.
    \label{eq:derivative_covariance}
\end{equation}

The constrained slow-feature objective can be written as
\begin{equation}
    J(U) = \operatorname{tr}(U^\top \Sigma_{\dot{x}} U)
    \quad
    \text{subject to}
    \quad
    U^\top \Sigma_x U = I_r.
    \label{eq:sfa_trace_objective}
\end{equation}
This is a generalized Rayleigh-Ritz problem. Introduce whitened coordinates
\begin{equation}
    z_t = \Sigma_x^{-1/2} x_t,
    \qquad
    V = \Sigma_x^{1/2} U.
\end{equation}
The constraint becomes
\begin{equation}
    V^\top V = I_r,
\end{equation}
and the objective becomes
\begin{equation}
    J(V)
    =
    \operatorname{tr}\left(
        V^\top
        \Sigma_x^{-1/2}\Sigma_{\dot{x}}\Sigma_x^{-1/2}
        V
    \right).
\end{equation}
Define the whitened derivative covariance
\begin{equation}
    M = \Sigma_x^{-1/2}\Sigma_{\dot{x}}\Sigma_x^{-1/2}.
\end{equation}
By the Rayleigh-Ritz theorem, the global minimizers are the $r$ eigenvectors of
$M$ with smallest eigenvalues, up to an orthogonal rotation when eigenvalues
are tied.

We now relate these eigenvectors to the slow eigenspace of $A$. Assume first
the clean case where $A$ is normal in the $\Sigma_x$-whitened inner product and
where $Q$ is isotropic after whitening. Then there exists an orthonormal basis
$Q_A=[q_1,\ldots,q_d]$ such that
\begin{equation}
    A q_i = \lambda_i q_i.
\end{equation}
Equivalently, after ordering modes from slow to fast,
\begin{equation}
    Q_A^{-1} A Q_A
    =
    \begin{bmatrix}
        \Lambda_s & 0 \\
        0 & \Lambda_f
    \end{bmatrix},
    \qquad
    \tilde{x}_t = Q_A^{-1}x_t
    =
    \begin{bmatrix}
        x_t^{(s)} \\
        x_t^{(f)}
    \end{bmatrix},
\end{equation}
where $\Lambda_s=\diag(\lambda_1,\ldots,\lambda_r)$ contains the slow modes
and $\Lambda_f$ contains the remaining fast modes. In these coordinates,
\begin{equation}
    \dot{\tilde{x}}_t
    =
    \begin{bmatrix}
        \Lambda_s-I & 0 \\
        0 & \Lambda_f-I
    \end{bmatrix}
    \tilde{x}_t
    +
    \tilde{w}_t.
\end{equation}
In this basis, the deterministic contribution to derivative variance along
mode $q_i$ is
\begin{equation}
    \|(A-I)q_i\|_2^2
    =
    |\lambda_i-1|^2.
\end{equation}
With isotropic whitened process noise, the noise contribution adds the same
constant to every mode. Thus the derivative covariance eigenvalue associated
with $q_i$ has the form
\begin{equation}
    \mu_i = |\lambda_i-1|^2 + \sigma_w^2.
    \label{eq:mode_derivative_variance}
\end{equation}
TODO: Replace this simplified isotropic-noise expression with the exact
generalized eigenvalue formula for non-isotropic $Q$.

The slow-fast spectral gap assumption states that there is an index set
$\mathcal{I}_s$ with $|\mathcal{I}_s|=r$ such that
\begin{equation}
    |\lambda_i-1|^2 < |\lambda_j-1|^2
    \qquad
    \text{for all } i \in \mathcal{I}_s,\; j \notin \mathcal{I}_s.
\end{equation}
Equivalently, the $r$ smallest derivative variances $\mu_i$ correspond exactly
to the slow modes. Hence the $r$-dimensional minimizing eigenspace of $M$ is
\begin{equation}
    \operatorname{span}\{q_i : i \in \mathcal{I}_s\}
    =
    \mathcal{S}_{\mathrm{true}}.
\end{equation}

Returning to the original coordinates, $U=\Sigma_x^{-1/2}V$. This change of
coordinates preserves the identified generalized eigenspace in the
$\Sigma_x$-inner product. Therefore any population minimizer $U_\star$ spans
the same slow subspace as the true slow latent component, up to an invertible
rotation within that subspace.

It remains to connect the population statement to the empirical objective. With
sufficient data from the stationary process, the empirical covariance
$\widehat{\Sigma}_x$ and empirical derivative covariance
$\widehat{\Sigma}_{\dot{x}}$ converge to their population values. Standard
perturbation results for eigenspaces, such as Davis-Kahan, imply that if the
spectral gap is nonzero then the empirical minimizing subspace converges to
the population slow subspace.

TODO: Add the explicit finite-sample perturbation bound and state whether the
desired result is consistency in probability or an asymptotic almost-sure
statement.
\end{proof}
